Dois pragmas úteis: "(hashtag)pragma GCC optimize("opt1,opt2...")" e "(hashtag)pragma GCC target("t1,t2...")".
Note que não tem espaço entre vírgulas.

Coloca isso no começo do código (antes dos includes). Lista de carinhas úteis:

\begin{itemize}
    \item \textbf{optimize}: O3, Ofast (toma cuidado com math com esse), unroll-loops (ja ta no O3 e Ofast), 
    tree-vectorize (precisa disso para usar avx, se não tiver O(hashtag) ou Ofast ativado)
    \item \textbf{target}: avx2 (se nao der, tenta avx), popcnt (otimiza familia do popcount), 
    lzcnt (otimiza familia do clz), abm, bmi, bmi2 (bmi não é subset do bmi2).
\end{itemize}

Se quiser ativar a otimização só em uma função, dá para por:
\begin{verbatim}
    __attribute__((target("t1,t2..."), optimize("opt1,opt2..."))) void funcao() ....
    \end{verbatim}